%!TEX root = ../LLMSeminar.tex
% ──────────────────────────────────────────────────────────────
%  Section 4 — Beyond the Loop
% ──────────────────────────────────────────────────────────────

\section{Beyond the loop}\label{sec:beyond}

\subsection{The real workflow}

You do not sit and watch the agent type.  You launch it, walk away, and
check the result.  The simplest orchestrator is a bash loop:

\begin{lstlisting}
while true; do
    claude --prompt "Fix the failing tests" \
           --output result.json
    if tests_pass result.json; then
        break
    fi
    sleep 10
done
\end{lstlisting}

Infrastructure at its most essential is a \code{while} loop in bash.

\subsection{Multi-agent orchestration}

The natural extension of the single-agent loop is \textbf{multi-agent
orchestration}: multiple LLM instances, each with a specialised role,
coordinated by a supervisor.

\begin{center}
\begin{tikzpicture}[node distance=1.8cm, every node/.style={align=center}]
  \node[sysbox, fill=spacecadet, text=white,
        minimum width=4.5cm, minimum height=0.9cm, font=\small]
    (sup) {\textbf{Supervisor}\\{\scriptsize orchestrates the pipeline}};

  \node[sysbox fill, below left=1.2cm and 0.8cm of sup,
        minimum width=3cm, font=\small]
    (a1) {\textbf{Agent 1}\\{\scriptsize writes code}};
  \node[sysbox fill, below=1.2cm of sup,
        minimum width=3cm, font=\small]
    (a2) {\textbf{Agent 2}\\{\scriptsize reviews code}};
  \node[sysbox fill, below right=1.2cm and 0.8cm of sup,
        minimum width=3cm, font=\small]
    (a3) {\textbf{Agent 3}\\{\scriptsize runs tests}};

  \draw[sysarrow] (sup) -- (a1);
  \draw[sysarrow] (sup) -- (a2);
  \draw[sysarrow] (sup) -- (a3);

  \draw[sysarrow, dashed, color=munsell] (a1.east) -- (a2.west);
  \draw[sysarrow, dashed, color=munsell] (a2.east) -- (a3.west);
\end{tikzpicture}
\end{center}

One agent writes code, another reviews it, a third runs tests.  The
supervisor decides when to iterate and when to stop.  This is useful
because agents are only moderately good at reasoning, but they are
\emph{much better at finding mistakes}.  By having a critique agent
review a proposal, you get a probabilistic oracle for correctness that
dramatically improves the output quality.

\begin{remark}
  Many frontier labs are reportedly running exactly this kind of
  multi-agent critique behind the scenes.  When the model appears to
  ``reason,'' part of that capability comes from automated
  self-critique during generation.
\end{remark}

\subsection{The natural endpoint of the stack}

The progression is clean:
\[
  \text{single function} \;\to\; \text{agent} \;\to\; \text{multi-agent}.
\]
Each layer adds orchestration on top of the same primitive:
$f : \Str \to \Str$.

\subsection{Limitations}

Every layer we built today has failure modes:

\begin{itemize}
  \item \textbf{Context rot}---older information degrades as the context
    fills.
  \item \textbf{Hallucination}---the model confidently produces false
    statements.
  \item \textbf{Tool misuse}---the agent may call tools incorrectly or
    in harmful ways.
  \item \textbf{Prompt fragility}---small changes in phrasing can
    produce wildly different results.
\end{itemize}

\begin{keyresult}[The human in the loop is not optional]
  The human in the loop is not optional.  LLMs have no fidelity.  They
  lie, they cheat, they steal, they flatter.  You need to give them
  fidelity as the expert, as the source of ground truth, through
  feedback.  The faster you can get your feedback loops, the better.
\end{keyresult}

\subsection{The automation spectrum}

Where do AI tools sit in the history of cognitive automation?

\begin{center}
\begin{tikzpicture}
  \draw[spacecadet, line width=1.5pt] (-5,0) -- (5,0);
  \node[above, font=\bfseries\small, text=spacecadet] at (-5,0.2) {Labour};
  \node[above, font=\bfseries\small, text=spacecadet] at (5,0.2) {Cognition};
  \fill[spacecadet] (-3.8,0) circle (3pt);
  \node[below=6pt, font=\small] at (-3.8,0) {Calculator};
  \fill[spacecadet] (-1.2,0) circle (3pt);
  \node[below=6pt, font=\small] at (-1.2,0) {Computer algebra};
  \fill[spacecadet] (1.2,0) circle (3pt);
  \node[below=6pt, font=\small] at (1.2,0) {Automated proving};
  \fill[cgblue] (3.8,0) circle (3pt);
  \node[below=6pt, font=\small, text=spacecadet] at (3.8,0)
    {\textbf{AI tools}};
\end{tikzpicture}
\end{center}

Every previous step on this spectrum was met with resistance and then
became indispensable.  AI tools are the next step.
